% Pitch marks above syllables.  A rising pitch is marked with one or more
% forward slashes "/" -- an "inverted backslash" -- and a falling pitch with
% backslashes "\".  Both are ordinary characters set in \markfont, so -- unlike
% the \pdfliteral strokes used before -- nothing here depends on the output
% driver, and the file is fully portable.  \risetwo/\risethree and
% \falltwo/\fallthree use two and three strokes.  \hold{...} (see hold.tex)
% underlines a held syllable; since it sits below the baseline it composes with
% the marks above, e.g. \rise{\hold{heer}}.
%
% \markfont sizes the marks (its slash/backslash glyphs are the strokes); scale
% it to make the marks bigger or smaller.  The backslash lives at code 92 in the
% 8t (T1) encoding of phvr8t.
\font\markfont=phvr8t scaled 1200   % marks: Helvetica; "/" = rise, "\" = fall
\chardef\bslash=92                  % backslash character code
% \riseheight is the fixed distance from the text baseline to a mark's baseline,
% so every mark sits the same distance above the text and they line up
% horizontally (rather than following each syllable's own height).  Its value is
% measured from the body font in layout.tex; reset it (e.g. per paragraph) to
% move all following marks together.
\newdimen\riseheight
% \risemark{strokes}{syllable} is the shared core: it measures the syllable's
% width, then \rlap overlays the stroke characters -- horizontally centered over
% the syllable and raised so their baseline sits \riseheight above the text
% baseline -- without consuming horizontal space, before printing the syllable
% itself.
\def\risemark#1#2{%
  \leavevmode                 % ensure horizontal mode, so a mark at the start
                              % of a line stays inline with the text
  \setbox0=\hbox{#2}%
  \rlap{\raise\riseheight\hbox to\wd0{\hss\markfont#1\hss}}%
  % No strut needed: room for the marks is provided globally by the fixed
  % \baselineskip set with the font, so all lines keep the same height.
  \box0
}
\def\rise{\risemark{/}}
\def\risetwo{\risemark{//}}
\def\risethree{\risemark{///}}
% \fall, \falltwo, \fallthree: the same, with backslashes for the falling pitch.
\def\fall{\risemark{\bslash}}
\def\falltwo{\risemark{\bslash\bslash}}
\def\fallthree{\risemark{\bslash\bslash\bslash}}
